\documentclass[letterpaper,11pt]{moderncv}
\moderncvstyle{banking}
\moderncvcolor{black}

\usepackage[letterpaper, margin=0.8in]{geometry}
\usepackage{xcolor}
\usepackage{fontspec}
\setmainfont{Times New Roman}
\usepackage[unicode]{hyperref}

\name{Saeyoung}{Kim}
\address{Cambridge, MA, USA}{}{}
\phone[mobile]{+1 (351) 322 5510}
\email{saeyoung\_kim@uml.edu}
\social[linkedin]{saeyoung-macx-kim}
\social[github]{PrimaryAnomaly}

\recipient{Hiring Committee}{Department of Electrical Engineering and Computer Sciences\\University of California, Berkeley\\Berkeley, CA}
\date{\today}
\opening{Dear Hiring Committee,}
\closing{Sincerely,}

\begin{document}

\makelettertitle

I am writing to apply for the Associate Development Engineer (4758C) position in EECS at UC Berkeley. I am a hands-on hardware-software integration engineer with experience building, integrating, testing, and troubleshooting embedded and electromechanical systems across academic research and industry environments. The role's emphasis on instructional lab support, embedded platforms, instrumentation, and dependable day-to-day technical operations closely matches my background.

In my current postdoctoral role at the University of Massachusetts Lowell, I develop Python tools, data pipelines, and model-based workflows for pharmaceutical manufacturing research. This work has strengthened my ability to translate experimental needs into practical technical systems, support multidisciplinary users, document methods clearly, and maintain reliable tools for repeated use in research settings.

Previously, at Hanwha Systems, I supported military SAR satellite programs as a Systems Integration Engineer. I developed and executed test procedures, defined integration and validation plans, contributed to Electrical Ground Support Equipment architecture, and diagnosed hardware-software interaction issues across networked subsystems. During my PhD in Electrical Engineering, I designed and fabricated multi-modal sensor systems, developed MCU firmware in C/C++, built embedded prototypes from circuit design through PCB layout and assembly, and debugged integrated electromechanical systems for autonomous operation.

I am especially interested in this opportunity because it combines hands-on engineering with direct support for learning and experimentation. I would be excited to contribute practical support across electronics, embedded systems, instrumentation, and experimental setup development while helping maintain lab spaces that are safe, dependable, and effective for students, faculty, and staff.

Thank you for your time and consideration. I would welcome the opportunity to discuss how my background in embedded systems, instrumentation, integration testing, and research support can contribute to the EECS instructional labs at Berkeley.

\makeletterclosing

\end{document}
